% !TeX program = xelatex
% UTF-8 transcription of all 19 source photographs, in filename/page order.
% Original wording, abbreviations, doses, and handwritten annotations retained.
% pic01.jpg = 20261012_192842739.jpg
% pic02.jpg = 20261012_192850431.jpg
% pic03.jpg = 20261012_192855248.jpg
% pic04.jpg = 20261012_192901433.jpg
% pic05.jpg = 20261012_192907779.jpg
% pic06.jpg = 20261012_192915370.jpg
% pic07.jpg = 20261012_192921407.jpg
% pic08.jpg = 20261012_192927565.jpg
% pic09.jpg = 20261012_192932921.jpg
% pic10.jpg = 20261012_192940382.jpg
% pic11.jpg = 20261012_192947401.jpg
% pic12.jpg = 20261012_192954496.jpg
% pic13.jpg = 20261012_193000410.jpg
% pic14.jpg = 20261012_193006804.jpg
% pic15.jpg = 20261012_193013785.jpg
% pic16.jpg = 20261012_193027872.jpg
% pic17.jpg = 20261012_193035823.jpg
% pic18.jpg = 20261012_193044111.jpg
% pic19.jpg = 20261012_193051349.jpg
% Standalone document; may also be input into an existing XePersian document.
\ifdefined\settextfont
  \def\AHFStandalone{0}
\else
  \def\AHFStandalone{1}
  \documentclass[12pt,a4paper]{article}
  \usepackage[margin=2cm]{geometry}
  \usepackage{amsmath,amssymb,array,longtable}
  \usepackage{xepersian}
  \IfFontExistsTF{Amiri}{\settextfont{Amiri}}{\settextfont{FreeSerif}}
  \setlatintextfont{FreeSerif}
\fi
\providecommand{\en}[1]{\lr{#1}}
\newcommand{\AHFpage}[1]{\clearpage\section*{موضوع: براونوالد ـ ۴۹ ـ تشخیص و درمان \en{AHF}}{\footnotesize تاریخ: \en{/ /}\quad صفحهٔ #1\par}}
\newcommand{\AHFreading}[1]{[خوانش نامطمئن: #1]}
\ifnum\AHFStandalone=1 \begin{document}\fi

\AHFpage{۱}
\begin{itemize}
\item \en{AHF} یکی از شایع‌ترین علل بستری در افراد بالای ۶۵ سال است.
\item نمودار: \en{$AHF\longrightarrow$ New Onset}؛ شاخهٔ دیگر: \en{Decompensation}.
\item ۵۰\% بیماران \en{AHF} که بستری می‌شوند \en{EF} نرمال دارند و \en{HFpEF} هستند.
\item ریسک مرگ در بیمارستان برای \en{HFrEF} بیشتر است ولی ریسک بستری مجدد و پروگنوز طولانی‌مدت برای هر دو گروه یکی است.
\item سن متوسط بیماران \en{AHF}، ۷۵ سال است. شیوع کلی نیز \en{F=M} است.
\item بیماران \en{F} در سن بالاتری بستری می‌شوند و احتمال اینکه \en{HFpEF} باشد بالاتر است.
\item شایع‌ترین کوموربیدیتی، \en{HTN} است که در \en{$2/3$} بیماران دیده می‌شود.
\item \en{CAD} در \en{$1/2$} و دیس‌لیپیدمی در \en{$1/3$} بیماران.
\end{itemize}
\subsection*{ارزیابی بیمار \en{AHF}}
در برخورد با بیمار \en{AHF} باید چند اقدام اولیه انجام دهیم:
\begin{enumerate}
\item تأیید تشخیص \en{AHF}.
\item درمان وضعیت تهدیدکنندهٔ حیات (شوک، نارسایی تنفسی).
\end{enumerate}

\AHFpage{۲}
\begin{enumerate}\setcounter{enumi}{2}
\item تشخیص و درمان بیماری‌ها / شرایط مرتبط (\en{ACS}، \en{PE}).
\item ارزیابی ریسک بیماران برای تریاژ (\en{ICU}، بخش).
\item تعیین پروفایل بالینی بیمار برای شروع درمان (\en{BP}، وضعیت کلیه).
\end{enumerate}
\subsection*{علائم \en{AHF}}
\begin{itemize}
\item شایع‌ترین علائم بیماران \en{AHF}، علائم مربوط به \en{Congestion} هستند.
\item دیس‌پنه شایع‌ترین علامت است و در ۹۰\% بیماران وجود دارد.
\item علائم کانجسشن: وریدهای سیستمیک.
\item ادم محیطی، افزایش وزن، افزایش سایز دور شکم.
\item علائم آتیپیک (خصوصاً در بیماران مسن).
\item خستگی، آنورکسی، \en{MSC}، اختلال خواب.
\item \en{Bendopnea}: احساس دیس‌پنه وقتی خم می‌شوند.
\end{itemize}
\subsection*{دسته‌بندی علائم}
\begin{itemize}
\item ناشی از \en{Volume Overload}.
\item \en{Symptom}: دیس‌پنه، سرفه، ارتوپنه \en{$\longleftrightarrow$} \en{Sign}: رال، پلورال افیوژن.
\item احساس ناخوشایند اندام تحتانی \en{$\longleftarrow$} ادم محیطی.
\item دیسکامفورت شکمی، سیری \en{$\longleftarrow$} آسیت، افزایش دور شکم، درد \en{RUQ}، هپاتومگالی، اسپلنومگالی.
\item افزایش \en{JVP}.
\end{itemize}

\AHFpage{۳}
\subsection*{ناشی از هایپوپرفیوژن}
\begin{itemize}
\item خستگی \en{$\longleftarrow$} سردی انتهاها.
\item \en{MSC}، خواب‌آلودگی در طی روز، اختلال تمرکز، کنفیوژن \en{$\longleftarrow$} افت \en{BP}، رنگ‌پریدگی، تغییر رنگ پوست.
\item پره‌سنکوپ، سنکوپ \en{$\longleftarrow$} فشار نبض پایین، فشار وریدی بالا.
\end{itemize}
\subsection*{سایر یافته‌ها}
\begin{itemize}
\item آنورکسی \en{$\longleftarrow$} هایپوولمی (ارتواستاتیک هایپوتنشن).
\item اختلال خواب \en{$\longleftarrow S_4$}.
\item تپش قلب \en{$\longleftarrow$} سوفل‌های سیستولیک و دیاستولیک.
\end{itemize}
\subsection*{معاینهٔ فیزیکی}
\begin{itemize}
\item تعیین دو نکته:
  \begin{itemize}
  \item آیا پرفیوژن کافی است؟ \en{Cold vs Warm}.
  \item آیا کانجسشن دارد؟ \en{Wet vs Dry}.
  \end{itemize}
\item اندازه‌گیری فشار خون:
  \begin{itemize}
  \item ۵۰\% بیماران: \en{$140<SBP$}.
  \item ۱۲\% بیماران: \en{$180<SBP$}؛ \en{HTN} زمینه‌ای + کریز هایپرتنسیو.
  \item ۲\% بیماران: \en{$SBP<90$}.
  \end{itemize}
\item بیمار با هایپوپرفیوژن سیستمیک ممکن است \en{BP} نرمال داشته باشد و بیمار با پرفیوژن طبیعی ممکن است به صورت مزمن \en{BP} پایین داشته باشد.
\end{itemize}

\AHFpage{۴}
\begin{itemize}
\item به همین دلیل \en{PP} (فشار نبض، \en{SBP-DBP}) نیز برای ما مهم است.
\item یک \en{low PP} نشان‌دهندهٔ \en{low Cardiac Output} است.
\item \en{High PP} \en{$\longleftarrow$} تیروتوکسیکوز، آنمی، رگورژیتاسیون آئورت.
\item بررسی \en{JVP} نشان‌دهندهٔ فشار ورید سیستمیک است.
  \begin{itemize}
  \item نشان‌دهندهٔ فشار دهلیز راست است.
  \item اغلب بیانگر فشار پرشدگی \en{LV} است.
  \item در بیمار با \en{TR}، انفارکت \en{RV}، \en{PHTN} قابل اعتماد نیست.
  \end{itemize}
\item رال و کراکل‌های تنفسی در ۷۰–۹۰\% بیماران \en{AHF} وجود دارد.
  \begin{itemize}
  \item در بیماران با سابقهٔ \en{HF} مزمن و \en{PHTN} به دلیل افزایش فعالیت لنفاوی ممکن است وجود نداشته باشد.
  \item نبود رال به معنای فشار پرشدگی طبیعی \en{LV} نیست.
  \end{itemize}
\item انتهاها سرد در حالی که نبض پریفرال خوب است، به معنا:
  \begin{itemize}
  \item \en{Marginal Cardiac Index}.
  \item \en{Marked Vasoconstriction}.
  \item یا هر دو.
  \end{itemize}
\item ادم محیطی در ۶۵\% بیماران \en{AHF} وجود دارد.
\end{itemize}

\AHFpage{۵}
\begin{itemize}
\item ادم محیطی در بیمارانی که با شوک کاردیوژنیک یا \en{LVEF} پایین مراجعه می‌کنند کمتر دیده می‌شود.
\item ادم در \en{AHF}، قرینه، \en{pitting}، \en{dependent} است.
  \begin{itemize}
  \item باید حداقل \en{4 lit} مایع خارج سلولی جمع شده باشد تا ادم قابل تشخیص ببینیم.
  \end{itemize}
\end{itemize}
\subsection*{بیومارکرها در \en{AHF}}
\begin{itemize}
\item برای ارزیابی بیمار در \en{ED} با دیس‌پنه، \en{COR=I}، ارزیابی \en{BNP} و \en{NT-proBNP} می‌باشد.
\item ارزش تشخیصی منفی بالاتر از ارزش تشخیصی مثبت است.
  \begin{itemize}
  \item اگر منفی باشد \en{HF} علت دیس‌پنه نیست.
  \item مثبت کاذب: \en{MI}، \en{PE}.
  \item منفی کاذب: چاقی.
  \end{itemize}
\item نمی‌تواند به صورت دقیق \en{HFrEF} را از \en{HFpEF} افتراق دهد.
\item سطح تروپونین در بیمار \en{AHF} در \en{ED} معمولاً بالاست.
  \begin{itemize}
  \item افزایش آن با پروگنوز بدتر همراه است.
  \item اندازه‌گیری آن \en{COR=I} \en{$\longleftarrow$} هم پروگنوز هم \en{ACS} همزمان.
  \item مثبت شدن آن در \en{AHF} بدون شواهد بالینی \en{ACS} لزوماً به معنی \en{MI type II} نیست.
  \end{itemize}
\end{itemize}

\AHFpage{۶}
\subsection*{سایر بررسی‌های آزمایشگاهی}
\begin{itemize}
\item عملکرد کلیه \en{$\longleftarrow$} \en{BUN, Cr}، \en{eGFR}.
\item \en{BUN} بیشتر از \en{Cr} با شدت بیماری \en{AHF} ارتباط دارد.
\item اگر خواستیم سایر تشخیص‌ها را \en{R/O} کنیم:
  \begin{itemize}
  \item \en{D-dimer}: آمبولی ریوی.
  \item پروکلسیتونین: عفونت.
  \end{itemize}
\end{itemize}
\subsection*{گرافی قفسهٔ سینه}
در بدو ورود بیمار با دیس‌پنه باید گرفته شود.
\subsection*{\en{EKG}}
\begin{itemize}
\item در همهٔ بیماران با شک به \en{AHF} باید اخذ شود.
\item در \en{AHF} احتمال تروپونین بالا است، شواهد ایسکمی در \en{EKG} مهم است.
\item آریتمی‌ها می‌توانند شروع‌کنندهٔ \en{AHF} باشند.
\item ۲۰–۳۰\% بیماران \en{AHF} ریتم \en{AF} دارند. (واژهٔ «اکثر» در این سطر خط خورده است.)
\end{itemize}
\subsection*{اکوکاردیوگرافی}
\begin{itemize}
\item در همهٔ بیماران \en{AHF} باید گرفته شود.
\item \en{RWMA}، اختلال عملکرد سیستولیک و دیاستولیک، عملکرد دریچه‌ها.
\item \en{$15<E:E_a\ \text{ratio}$} مطرح‌کنندهٔ \en{$15\,mmHg<PCWP$}.
\end{itemize}

\AHFpage{۷}
درمان \en{AHF} دارای چهار فاز می‌باشد.
\subsection*{فاز \en{I}: مراقبت \en{Urgent/Emergent}}
\begin{itemize}
\item تشخیص قطعی \en{AHF}.
\item درمان شرایط تهدیدکنندهٔ حیات.
\item شروع درمان با هدف رفع علائم بالینی.
\item تشخیص علت دکامپنسیشن.
\item اگر \en{$SpO_2<90$}: \en{$O_2$ therapy}.
  \begin{itemize}
  \item باید مراقبت باشیم که \en{$0.4<FiO_2$} در بیمار با اختلال عملکرد سیستولیک می‌تواند \en{hyperoxia-induced Vasocons} ایجاد کند.
  \end{itemize}
\item اگر بیمار \en{COPD} است نباید تا حداکثر \en{$SpO_2$} را بالا ببریم (بالای ۹۴\%)، چرا که خطر سرکوب تنفسی دارد.
\item \en{CPAP} و \en{NIPPV} به بهبود بیماران \en{AHF} کمک می‌کنند؛ در نمودار هر دو زیر عنوان \en{NIV} قرار دارند.
\item کنتراندیکاسیون \en{NIV}: نیاز به انتوبیشن، هوشیار و \AHFreading{کوآپ نبودن بیمار}.
\item استفاده از مورفین با ریسک بالاتر نیاز به \en{MV}، \en{ICU}، مرگ همراه است.
\item همهٔ بیماران \en{AHF} نیاز به دیورتیک ندارند. اگر مشکل بیمار والیوم ری‌دیستریبیوشن باشد، به جای والیوم اورلود، به جای دیورتیک نیاز به وازودیلاتور دارند.
\end{itemize}

\AHFpage{۸}
\subsection*{اندیکاسیون درمان \en{urgent} با دیورتیک}
\begin{itemize}
\item وجود \en{Congestion} \en{$\longleftarrow$} دیلاتاسیون وریدی سیستمیک، هایپرتنشن ریوی.
\item ادم.
\item دوز دیورتیک به صورت بولوس: ۱–۲٫۵ برابر دوز فوروزماید خوراکی را به صورت وریدی می‌دهیم؛ \en{chronic User}. (عدد قبلی خط خورده و ۱–۲٫۵ بالای آن نوشته شده است.)
\item اگر فشار پایین نباشد \en{$\longleftarrow$} وازودیلاتور می‌دهیم.
\end{itemize}
\subsection*{اندیکاسیون بستری \en{AHF}}
\begin{itemize}
\item \en{$BP\downarrow$}، \en{MSC}، \en{AKI}.
\item هایپوکسی، \en{ACS}، افزایش وزن \en{$5\,kg<$}.
\item وجود آریتمی که از نظر همودینامیکی مهم باشد (\en{AF+RVR}).
\item اختلال الکترولیتی.
\end{itemize}
\subsection*{فاز \en{II} (مراقبت در بیمارستان)}
\begin{itemize}
\item اندازه‌گیری روزانهٔ وزن.
\item چارت \en{I/O}.
\item چک علائم حیاتی (شامل فشار ارتواستاتیک).
\item اندازه‌گیری روزانهٔ علائم و نشانه‌ها.
\item چک روزانهٔ \en{BUN, Cr, Na, K}.
\end{itemize}

\AHFpage{۹}
\begin{itemize}
\item محدودیت سدیم روزانه کمتر از \en{2 gr}.
\item محدودیت مایعات روزانه کمتر از \en{2 lit}.
\item پروفیلاکسی \en{DVT} مگر کنتراندیکاسیون واقعی وجود داشته باشد.
\item در این فاز با توجه به احتمال اختلال عملکرد کلیه ممکن است نیاز باشد داروهای سیستم \en{RAAS} (\en{ACE-i, ARB, ARNI}) را قطع یا تعدیل دوز کنیم.
  \begin{itemize}
  \item به طور کلی توصیه می‌شود درمان‌های براساس گایدلاین تا حد امکان قطع نشوند.
  \end{itemize}
\item اگر بیمار \en{BB} مصرف می‌کند نباید آن را قطع نماییم (حین \en{AHF}) مگر آنکه افت فشار شدید یا شوک کاردیوژنیک وجود داشته باشد.
\item این فاز بهترین زمان تبدیل داروی بیمار به ساکوبیتریل ـ والسارتان است.
\end{itemize}
\subsection*{سندرم کاردیو ـ رنال در بیماران بستری}
\begin{itemize}
\item درمان \en{Volume Overload} در بیمار \en{AHF} به درمان مقاوم است یا راجعه می‌باشد، به دلیل نارسایی کلیهٔ پیش‌رونده.
\item علی‌رغم شواهد بالینی یا همودینامیک \en{Congestion}، عدد \en{Cr} بیش از \en{0.3} بالا می‌رود یا \en{GFR} بیش از ۲۵\% کاهش می‌یابد. (حاشیه: \en{GFR} مهم است؛ دور آن دایره کشیده شده است.)
\item در ۲۵–۳۰\% بیمارانی که با \en{AHF} بستری می‌شوند وجود دارد.
\end{itemize}

\AHFpage{۱۰}
\begin{itemize}
\item تشخیص سندرم کاردیو ـ رنال آسان ولی منج آن سخت است.
\item اگر \en{$80<BUN$} و \en{$3<Cr$} یا هایپرکالمی:
  \begin{itemize}
  \item قطع داروهای \en{RAAS}، در صورت نیاز وازودیلاتورهای دیگر می‌توانند داده شوند.
  \item وریدی: نیتروگلیسیرین، نیتروپروساید.
  \item خوراکی: هیدرالازین، ایزوسورباید دی‌نیترات.
  \end{itemize}
\item اغلب باید دوز دیورتیک را بالا برد \en{$\longleftarrow$} ریسک مقاومت وجود دارد.
\item سطح سدیم ادرار معیاری مناسب برای ارزیابی پاسخ به دیورتیک است.
\end{itemize}
\subsection*{تشدید \en{HF} حین بستری (\en{WHF})}
\begin{itemize}
\item باید بدانیم که بیماران مختلف سیر بالینی بسیار متفاوتی برای \en{AHF} دارند.
\item بدتر شدن شرایط بالینی که نیازمند تشدید درمان است.
  \begin{itemize}
  \item علائم و نشانه‌های \en{HF} تشدید می‌شوند.
  \end{itemize}
\end{itemize}
\subsection*{فاز \en{III} (برنامه‌ریزی برای ترخیص)}
\begin{itemize}
\item بیمار نباید شواهد بالینی \en{Congestion} داشته باشد.
\item باید بتواند پله را بالا برود.
\item داروهایی که \en{Outcome} بیمار \en{HF} را بهبود می‌بخشند باید شروع شوند.
  \begin{itemize}
  \item اگر از قبل مصرف می‌شده نباید آن‌ها را قطع کرد مگر کنتراندیکاسیون جدید داشته باشد.
  \end{itemize}
\end{itemize}

\AHFpage{۱۱}
\subsection*{فاز \en{IV} (مراقبت پس از ترخیص)}
ویزیت ۷–۱۰ روز پس از ترخیص از بیمارستان.
\subsection*{الگوریتم برخورد اولیه با بیمار \en{AHF}}
\begin{latin}
\textbf{Urgent phase}
\begin{enumerate}
\item Cardiogenic Shock?
  \begin{itemize}
  \item Yes $\longrightarrow$ Circulatory Support: Pharmacologic; Mechanical $\longrightarrow$ Transfer to ICU/CCU.
  \item No $\longrightarrow$ Respiratory Failure?
  \end{itemize}
\item Respiratory Failure?
  \begin{itemize}
  \item Yes $\longrightarrow$ Ventilatory Support: $O_2$; NIV; MV.
  \item No $\longrightarrow$ Immediate phase.
  \end{itemize}
\end{enumerate}
\textbf{Immediate phase: 60--120 min}
\par Identification of acute Etiology (also reached after Transfer to ICU/CCU):
\begin{center}
\begin{tabular}{ll}
C & ACS \\
H & HTN Emergency \\
A & Arrythmia \\
M & Acute Mechanical Cause \\
P & PE \\
\end{tabular}
\end{center}
\begin{itemize}
\item Yes $\longrightarrow$ Immediate Specific Therapy.
\item No $\longrightarrow$ Diagnostic Workup.
\end{itemize}
\end{latin}

\AHFpage{۱۲}
\subsection*{الگوریتم درمان \en{AHF} در اورژانس بر اساس وجود \en{Congestion}}
\begin{latin}
\textbf{Assessment of Hemodynamic Profile}
\par Presence of Congestion?
\begin{itemize}
\item Yes (95\% of pts) $\longrightarrow$ Wet patient.
\item No (5\% of pts) $\longrightarrow$ Dry patient.
\end{itemize}
\par Adequate Peripheral Perfusion?
\begin{itemize}
\item Dry patient, Yes $\longrightarrow$ Dry and Warm = Compensated $\longrightarrow$ Adjust oral Therapy.
\item Dry patient, No $\longrightarrow$ Dry and Cold; hypoperfused; hypovolemic $\longrightarrow$ Fluid challenge; Inotropic agent.
\item Wet patient, Yes ($BP$ Nl or $\uparrow$) $\longrightarrow$ Wet and Warm:
  \begin{itemize}
  \item Vascular Type: fluid Redistribution, HTN Predominates $\longrightarrow$ Vasodilator; Diuretic.
  \item Cardiac Type: fluid accumulation, Congestion Predominates $\longrightarrow$ Diuretic; Vasodilator; Ultrafiltration (if diuretic resistance).
  \end{itemize}
\item Wet patient, No $\longrightarrow$ Wet and Cold pts $\longrightarrow SBP<90\,mmHg$?
  \begin{itemize}
  \item Yes: Inotropic Agents; Vasopressor; Diuretic (with perfusion corrected); Mechanical Circulatory Support.
  \item No: Vasodilator; Diuretic; Inotropic Agents.
  \end{itemize}
\end{itemize}
\end{latin}

\AHFpage{۱۳}
\subsection*{اپروچ درمانی برای مدیریت حجم در بیماران \en{AHF}}
شدت \en{Overload}:
\begin{itemize}
\item متوسط \en{$\longleftarrow$} فوروزماید، تا ۲٫۵ برابر دوز خوراکی یا \en{20--40 mg}، \en{IV} داده شود.
\item شدید \en{$\longleftarrow$} فوروزماید، تا ۲٫۵ برابر دوز خوراکی یا \en{40--120 mg}، بولوس \en{IV}.
\item \en{Maintenance}: \en{5--40 mg/h inf}.
\item مقاوم به دیورتیک لوپ \en{$\longleftarrow$} اضافه کردن:
  \begin{itemize}
  \item \en{HCTZ} (دور نام دایره کشیده شده): \en{25--50 mg BD}.
  \item یا متولازون: \en{2.5--10 mg daily}.
  \item یا اسپیرونولاکتون: \en{25--50 mg daily}؛ \en{if No AKI and Nl or low Serum $K^+$}.
  \end{itemize}
\item اگر آلکالوز وجود داشت \en{$\longleftarrow$} استازولاماید \en{$\longleftarrow$} \en{0.5 mg iv}.
\item مقاوم به لوپ و تیازید \en{$\longleftarrow$} الترافیلتریشن یا همودیالیز.
\end{itemize}
\subsection*{برخورد با بیمار \en{AHF + AF, RVR}}
\begin{itemize}
\item شایع‌ترین تاکی‌آریتمی پایدار در بیماران \en{AHF} است.
\item قابل افتراق نیست که \en{AF} نتیجهٔ \en{AHF} بوده یا علت آن.
\item اگرچه که پاسخ بطنی سریع به دنبال بهبود دیس‌پنه بهتر می‌شود ولی گاهی وقت‌ها نیاز به درمان هست.
\end{itemize}

\AHFpage{۱۴}
\begin{itemize}
\item کاردیوورژن اورژانسی معمولاً اندیکاسیون ندارد مگر در بیمار \en{Unstable}.
  \begin{itemize}
  \item ریسک عود \en{AF} خیلی بالاست.
  \end{itemize}
\item درمان در بیمار با اختلال سیستولیک:
  \begin{itemize}
  \item \en{BB} با احتیاط فراوان.
  \item دیگوکسین وریدی \en{$\longleftarrow$} نباید مسیر فرعی وجود داشته باشد.
  \item آمیودارون.
  \end{itemize}
\item \en{L-type CCB} (دیلتیازم) و سایر داروهای با عملکرد اینوتروپ منفی ممنوع هستند مگر \en{EF} نرمال باشد.
\end{itemize}
\subsection*{برخورد با بیمار \en{AHF} + نارسایی بطن راست}
\begin{itemize}
\item شایع‌ترین علت نارسایی \en{RV} در بیماران \en{AHF} نارسایی \en{LV} است.
\item فرم ایزولهٔ نارسایی \en{RV} بسیار نادر است ولی کشنده است:
  \begin{itemize}
  \item انفارکت حاد \en{RV}.
  \item \en{PE} حاد.
  \item \en{Severe pHTN}.
  \end{itemize}
\item اگر در زمینهٔ انفارکت باشد، بهترین اقدام ری‌پرفیوژن سریع است.
\item \en{PE} که همودینامیک را مختل کند \en{$\longleftarrow$} ترومبولیتیک.
\item هدف \en{CVP} در این بیماران \en{10--15 cm$H_2O$} است.
  \begin{itemize}
  \item در بیمارانی که بخواهیم مایع دادن را مانیتور کنیم.
  \end{itemize}
\end{itemize}

\AHFpage{۱۵}
\begin{itemize}
\item وازودیلاتورهای سلکتیو شریان‌های ریوی با کاهش افترلود عملکرد \en{RV} را بهبود می‌بخشند.
  \begin{itemize}
  \item استنشاقی: نیتریک‌اکساید، آنالوگ‌های پروستاسیکلین.
  \item وریدی: سیلدنافیل، آنالوگ‌های پروستاسیکلین.
  \end{itemize}
\item اگر بیمار تحت \en{MV} قرار گرفت:
  \begin{itemize}
  \item هدف \en{Normoxia} و \en{hypocarbia} است.
  \item \en{TV} متوسط \en{$\longleftarrow$} \en{8 ml/kg}.
  \item حداقل \en{PEEP} ممکن (\en{$<12\,cmH_2O$}).
  \end{itemize}
\end{itemize}
\subsection*{برخورد با بیمار \en{AHF + ACS}}
\begin{itemize}
\item افتراق بالینی بسیار دشوار است چون هم ممکن است یافته‌های \en{EKG} گول‌زننده باشد، هم بیومارکرها مثل \en{Trop}.
\item باید به طور تهاجمی \en{ACS} را درمان کنیم.
\item در نبود شوک کاردیوژنیک از \en{Inodilators} به هیچ عنوان استفاده نکنیم: \en{Milrinone, Levosimendan}.
\end{itemize}
\subsection*{برخورد با بیمار \en{AHF} + شوک کاردیوژنیک}
تعریف شوک کاردیوژنیک:
\begin{itemize}
\item افت فشار شدید (\en{$SBP<80\,mmHg$}).
\item بیشتر از \en{30 min} طول بکشد.
\item با کاهش شدید ایندکس قلبی همراه باشد: \en{$<1.8\,lit/min/m^2$}.
\item در حالی که فشار پرشدگی بطن چپ کافی است (\en{$PCWP\geq18\,mmHg$}).
\item منجر به \en{Organ Failure} شود.
\end{itemize}

\AHFpage{۱۶}
\begin{itemize}
\item کمتر از ۴\% بیماران \en{AHF} با شوک کاردیوژنیک مواجه هستند که معمولاً در زمینهٔ \en{MI} می‌باشد.
\item درمان: اینوتروپ‌های تزریقی وریدی؛ یا وازوکانستریکتورها.
\end{itemize}
\subsection*{برخورد با \en{BP} در بیمار با \en{AHF}}
\begin{itemize}
\item یکی از مهم‌ترین معیارهای پروگنوستیک، \en{BP} بیمار می‌باشد.
\item اغلب بیماران \en{BP} بالا دارند پس به راحتی می‌توانیم وازودیلاتور بدهیم.
\item وازودیلاتورها درمان پایه‌ای \en{AHF} در این بیماران:
  \begin{itemize}
  \item \en{AHF} + ادم ریوی.
  \item \en{AHF} که \en{Cold} است ولی \en{$85\text{--}100<SBP$} می‌باشد.
  \end{itemize}
\item کاربرد وازودیلاتور در بیمار \en{AHF}:
  \begin{itemize}
  \item کاهش پره‌لود: از طریق رفع انقباض وریدها.
  \item کاهش افترلود: از طریق رفع انقباض شریان‌ها.
  \end{itemize}
\item درمان علت زمینه‌ای \en{$BP\downarrow$} ضروری است (\en{$SBP<85\text{--}90$}).
  \begin{itemize}
  \item \en{ACS}، \en{PE}، ندرتاً هایپوولمی.
  \end{itemize}
\item هایپوتنشن بدون علامت: بیمار شواهد \en{Congestion} ندارد و شواهدی از هایپوپرفیوژن کلیوی و محیطی نیز دیده نمی‌شود.
  \begin{itemize}
  \item نیاز به درمان اورژانسی ندارد.
  \end{itemize}
\end{itemize}

\AHFpage{۱۷}
\begin{itemize}
\item استفاده از اینوتروپ وریدی شانس مرگ داخل بیمارستان را ۲٫۵ برابر می‌کند.
\end{itemize}
\subsection*{وضعیت حجم بیماران \en{AHF}}
\begin{itemize}
\item بیماری که شواهد \en{Overload} مانند ادم، آسیت و … دارد \en{$\longleftarrow$} دیورتیک.
\item بیمار با شواهد بالینی \en{Congestion} حداقل \en{4--5 lit} مایع اضافه دارد.
  \begin{itemize}
  \item احتمالاً عدد باید حتی بالای \en{10 lit} باشد.
  \end{itemize}
\item وازودیلاتور وریدی دیس‌پنه را سریع‌تر از دیورتیک رفع می‌کند.
\end{itemize}
\subsection*{عملکرد کلیه در بیماران \en{AHF}}
\begin{itemize}
\item وقتی دیورتیک می‌دهیم باید این مؤلفه را چک کنیم:
  \begin{itemize}
  \item عملکرد کلیه، الکترولیت‌ها، وضعیت حجمی بیمار.
  \end{itemize}
\item \en{$2/3$} بیماران \en{AHF} حداقل یک نارسایی متوسط کلیه دارند.
\item افت \en{GFR} با مقاومت به دیورتیک همراه است و باید دوز بالاتر بدهیم.
\end{itemize}
\subsection*{وازودیلاتورهای مورد استفاده در \en{AHF}}
\begin{itemize}
\item نیتروگلیسیرین:
  \begin{itemize}
  \item دوز شروع: \en{20 $\mu$g/min}.
  \item دوز هدف: \en{40--400 $\mu$g/min}.
  \item در هر دوز: بیش از ۲۴ ساعت، تاکی‌فیلاکسی.
  \end{itemize}
\item نیتروپروساید:
  \begin{itemize}
  \item دوز شروع: \en{0.3 $\mu$g/kg/min}.
  \item دوز هدف: \en{0.3--5 $\mu$g/kg/min} (۵ بالای عدد خط‌خورده نوشته شده است).
  \item \en{$<4\,\mu g/kg/min$}.
  \item در \en{AMI} ندهیم.
  \end{itemize}
\end{itemize}

\AHFpage{۱۸}
\subsection*{اینوتروپ‌ها در درمان \en{AHF}}
\begin{itemize}
\item \en{EN}: \en{0.05--0.5 $\mu$g/kg/min}.
  \begin{itemize}
  \item ریسک آریتمی، تشدید هایپوپرفیوژن اندام‌ها، ایسکمی قلبی.
  \end{itemize}
\item \en{NEN}: \en{0.2--1 $\mu$g/kg/min}.
  \begin{itemize}
  \item ایسکمی قلبی، آریتمی، تشدید هایپوپرفیوژن اندام‌ها.
  \end{itemize}
\item دوپامین: \en{1--2 $\mu$g/kg/min} \en{$\longrightarrow$} \en{2--4 $\mu$g/kg/min}.
  \begin{itemize}
  \item عوارض مشابه، نیز خطر کمتر برای \en{End Organ}.
  \end{itemize}
\item دوبوتامین: \en{1--2 $\mu$g/kg/min} \en{$\longrightarrow$} \en{2--20 $\mu$g/kg/min}.
  \begin{itemize}
  \item خطر کمتر برای \en{End Organ}.
  \end{itemize}
\item تیتراسیون دوز، ۲ برابر کردن است.
\end{itemize}
\subsection*{دربارهٔ دیورتیک‌های لوپ (فوروزماید) در \en{AHF}}
\begin{itemize}
\item شروع اثر در فرم وریدی ظرف ۳–۴ دقیقه.
\item دوز شروع در بیماری که مصرف مزمن دارد + کلیهٔ نارس دارد:
  \begin{itemize}
  \item اختلال عملکرد کلیوی / \en{Severe Volume Overload}
  \item ۲٫۵ برابر دوز خوراکی بیمار به صورت وریدی داده می‌شود.
  \end{itemize}
\item اگر به دوز بولوس اولیه جواب نداد آن را ۲ برابر می‌کنیم. (دور ۲ برابر دایره کشیده شده است.)
\item اندیکاسیون درمان انفوزیون:
  \begin{itemize}
  \item مقاومت به دیورتیک.
  \item \en{Overload} شدید بیشتر از ۱۰–۱۵ لیتر.
  \end{itemize}
\end{itemize}

\AHFpage{۱۹}
\begin{itemize}
\item تفاوتی بین درمان \en{BD} بولوس با انفوزیون مداوم وجود ندارد.
  \begin{itemize}
  \item دقیق نمی‌توان اظهارنظر کرد.
  \end{itemize}
\item در موارد مقاومت به دیورتیک می‌توان دیورتیک تیازیدمانند (اثر سینرژیستیک) را به درمان اضافه کرد:
  \begin{itemize}
  \item \en{HCTZ} وریدی (\en{500--1000 mg})، یا
  \item متولازون خوراکی (\en{2.5--10 mg}).
  \end{itemize}
\item قبل از هر دوز فوروزماید.
\item باید مراقب افت فشار بود.
  \begin{itemize}
  \item عملکرد کلیه، الکترولیت‌ها.
  \end{itemize}
\item \en{NSAID} اثربخشی دیورتیک‌ها را کم می‌کند و باید از آن‌ها پرهیز شود.
\item اگر هایپرکالمی وجود داشت \en{$\longleftarrow$} دیورتیک پتاسیم اسپیرینگ.
\end{itemize}

% Printed subject/date fields and handwritten page numbers are retained above.
% The notebook emblem contains no text.
\ifnum\AHFStandalone=1
  \let\AHFEnd\enddocument
\else
  \let\AHFEnd\relax
\fi
\AHFEnd
